\documentclass[11pt]{article}
\usepackage[margin=1in]{geometry}
\usepackage{amsmath,amssymb,amsthm}
\usepackage[colorlinks=true,urlcolor=blue,linkcolor=blue]{hyperref}
\newtheorem{theorem}{Theorem}
\title{A three-vertex counterexample to the\protect\\cospectral quiver mutation conjecture}
\author{Principia Math in Progress}
\date{28 September 2026}
\begin{document}
\maketitle

\section*{Task 1: resolve MathDB 335493}
For a cluster quiver $Q$, let $\operatorname{Mut}(Q)$ denote its mutation
class. Li and Liu's Conjecture 4.6 asks whether any two quivers in this
class with the same exchange spectrum can be connected by mutations
at sinks and sources, after allowing an initial quiver isomorphism
\cite{LL}. We disprove this assertion.

We use the convention $b_{ij}>0$ for $b_{ij}$ arrows from $i$ to $j$.
Changing the exchange-matrix sign convention does not change the proof.
Multiple arrows are permitted; loops and oriented two-cycles are not.
A sink has no outgoing arrows, and a source has no incoming arrows.

\begin{theorem}
Let $Q$ have three vertices and exactly the following arrow bundles:
\[
1\xrightarrow{\ 3\ }2,\qquad
2\xrightarrow{\ 4\ }3,\qquad
3\xrightarrow{\ 6\ }1.
\]
Then $Q$ and its opposite quiver $Q^{\mathrm{op}}$ belong to the same
mutation class and have equal exchange spectra. There is no quiver
isomorphic to $Q$ from which $Q^{\mathrm{op}}$ can be obtained by
mutations at sinks and sources.
\end{theorem}

\begin{proof}
Mutate at vertex $2$. The three arrows $1\to2$ and four arrows
$2\to3$ create twelve arrows $1\to3$. Cancelling against the six
existing arrows $3\to1$ leaves six arrows $1\to3$. Reversing the
arrows incident to vertex $2$ leaves three arrows $2\to1$ and four
arrows $3\to2$. Thus
\[
\mu_2(Q)=Q^{\mathrm{op}},\qquad
B(Q)=\begin{pmatrix}0&3&-6\\-3&0&4\\6&-4&0\end{pmatrix},
\qquad B(Q^{\mathrm{op}})=-B(Q).
\]
In particular both quivers lie in $\operatorname{Mut}(Q)$.
Direct expansion gives
\[
\det(\lambda I-B(Q))
=\lambda^3+(3^2+4^2+6^2)\lambda
=\lambda(\lambda^2+61)
=\det(\lambda I-B(Q^{\mathrm{op}})).
\]
Both exchange spectra are therefore $\{0,i\sqrt{61},-i\sqrt{61}\}$,
with the displayed algebraic multiplicities.

The quivers are not isomorphic. An isomorphism must preserve the number
of arrows between each unordered pair of vertices. The three such
numbers are distinct, so each of the three underlying edges must map
to itself. Intersecting two edges forces each vertex to be fixed.
The resulting identity vertex map cannot reverse an arrow, whereas
every arrow of $Q^{\mathrm{op}}$ is reversed from $Q$.

Finally, every vertex of $Q$ has both an incoming and an outgoing
arrow. This property is invariant under quiver isomorphism. Hence
no quiver $R$ isomorphic to $Q$ admits even a first mutation at a
sink or source. The only possible such sequence is the empty one;
it would require $R=Q^{\mathrm{op}}$, contradicting the
nonisomorphism just proved. This rules out every possible $R$ in
the conjecture and proves the counterexample.
\end{proof}

\paragraph{Source scope.}
The conjecture is unchanged between the 19-page arXiv manuscript
(Conjecture 4.6, p.~15) and the final 29-page article
(Conjecture 4.6, printed p.~1555). Both entire versions were read.
The neighboring Theorem 4.5 excludes three-cycles; the conjecture
has no such restriction. The source's definition allows multiple
arrows, so the example satisfies the intended hypotheses. The
argument does not address a modified conjecture restricted to
acyclic or finite-type mutation classes.

\paragraph{Verification and review.}
The accompanying integer-arithmetic verifier checks the mutation
formula, both characteristic polynomials, all six vertex permutations,
and the absence of sinks and sources. These checks certify the finite
example; the proof above explains why no arbitrarily long allowed
sequence can work. The live MathDB and bounded literature checks
found no prior full resolution claim. Their sources and limitations
are recorded in the separate audit. This is author-agent review,
not independent peer review.

\section*{Tasks 2--5: proposed follow-up problems}
These are separate research proposals, not completed results.
\begin{enumerate}
\setcounter{enumi}{1}
\item For each mutation class on three vertices, classify the
isomorphism classes of quivers having a prescribed exchange polynomial,
and determine their components under sink and source mutations.
\item Determine whether Conjecture 4.6 holds for every quiver whose
mutation class is of finite Dynkin type, including types $D$ and $E$.
\item Characterize all four-vertex quivers and vertices $k$ for which
mutation at $k$ preserves the exchange polynomial, using explicit
conditions on the six signed arrow multiplicities.
\item Decide whether two cospectral quivers in one mutation class must
be connected, up to isomorphism, by a sequence of mutations each of
which preserves the exchange polynomial; if false, construct a
counterexample of minimum vertex number.
\end{enumerate}

\begin{thebibliography}{9}
\bibitem{LL} Fang Li and Siyang Liu,
\emph{On exchange spectra of valued quivers and cluster algebras},
Pure Appl. Math. Q. \textbf{16} (2020), no.~5, 1533--1561.
\href{https://doi.org/10.4310/PAMQ.2020.v16.n5.a6}{doi:10.4310/PAMQ.2020.v16.n5.a6}.
Preprint: \href{https://arxiv.org/abs/1703.02225}{arXiv:1703.02225v1},
\emph{On Exchange Spectra of Valued Cluster Quivers and Cluster Algebras}.
\end{thebibliography}
\end{document}
